\chapter{Divisibilidad y algoritmo de Euclides}
\label{md:division}

Se usa $\N=\{1,2,\ldots\}$ y $\N_0=\N\cup\{0\}$.
El \textbf{principio de buena ordenación} afirma que todo subconjunto
no vacío de $\N_0$ tiene un mínimo. En $\Z$ la suma y el producto son
conmutativos y el producto distribuye sobre la suma; no hay divisores
de cero. Por ello, $ab=ac$ con $a\ne0$ implica $b=c$.

\section{Divisibilidad y división con resto}
\label{md:divisibilidad}

Para $a,b\in\Z$, con $a\ne0$, se escribe $a\mid b$ si existe $k\in\Z$
tal que $b=ak$. Así, $a$ es divisor de $b$ y $b$ es múltiplo de $a$.
Todo entero no nulo divide a $0$; $1$ y $-1$ dividen a cualquier entero.
El signo de un divisor no afecta a la divisibilidad.
Las propiedades básicas son
\[
 \begin{gathered}
 a\mid b,\ b\mid c\ \imp\ a\mid c,\qquad
 a\mid b,\ a\mid c\ \imp\ a\mid(bx+cy)\quad(x,y\in\Z),\\
 a\mid b,\ b\mid a\ \imp\ |a|=|b|.
 \end{gathered}
\]
Si $a,b>0$ y $a\mid b$, entonces $a\le b$.
El valor absoluto satisface $|ab|=|a||b|$ y
$|a+b|\le |a|+|b|$.

La \textbf{división euclídea} de $a\in\Z$ por $b\ne0$ determina
únicos $q,r\in\Z$ con
\[
 a=bq+r,\qquad 0\le r<|b|.
\]
Se llaman cociente y resto. La existencia se obtiene eligiendo el menor
entero no negativo de la forma $a-bq$; dos restos posibles no pueden
diferir en un múltiplo no nulo de $b$ dentro de ese intervalo.

Para $m>0$, $\MOD(a,m)$ designa este resto. Se puede reducir antes
o después de sumar o multiplicar:
\[
 \begin{aligned}
 \MOD(a+b,m)&=\MOD(\MOD(a,m)+\MOD(b,m),m),\\
 \MOD(ab,m)&=\MOD(\MOD(a,m)\MOD(b,m),m).
 \end{aligned}
\]
\begin{ejemplo}
$-17=5(-4)+3$. El resto es $3$, no $-2$: la división euclídea
exige un resto no negativo.
\end{ejemplo}

\section{Máximo común divisor y Bézout}
\label{md:bezout}

Si $(a,b)\ne(0,0)$, el \textbf{máximo común divisor}
$d=\mcd(a,b)$ es el mayor divisor común positivo. Además,
todo divisor común divide a $d$. Se tiene
\[
 \mcd(a,0)=|a|\quad(a\ne0),\qquad
 \mcd(a,b)=\mcd(|a|,|b|).
\]
Por convenio, $\mcd(0,0)=0$. Dos enteros son \textbf{primos entre sí}
o coprimos si su máximo común divisor es $1$.

La \textbf{identidad de Bézout} garantiza enteros $s,t$ tales que
\[
 \boxed{\mcd(a,b)=as+bt.}
\]
Su idea es tomar la menor combinación lineal positiva de $a,b$:
la división con resto muestra que divide a ambos, y cualquier divisor
común divide a esa combinación. En particular,
\[
 \mcd(a,b)=1\ \eqv\ \exists s,t\in\Z:\ as+bt=1.
\]
Si $d=\mcd(a,b)>0$, entonces $\mcd(a/d,b/d)=1$.
Para $k\ne0$,
$\mcd(ka,kb)=|k|\mcd(a,b)$.

\section{Cálculo por Euclides}
\label{md:euclides}

Si $a=bq+r$ y $b\ne0$, los divisores comunes de $a,b$ coinciden
con los de $b,r$. Por tanto,
\[
 \mcd(a,b)=\mcd(b,r).
\]
Para calcularlo se trabaja con valores absolutos y se encadenan divisiones:
\begin{enumerate}
 \item Dividir el mayor por el menor.
 \item Sustituir el par por divisor y resto; repetir mientras el resto no sea cero.
 \item El divisor de la división que da resto cero es el máximo común divisor.
\end{enumerate}
Los restos positivos decrecen, lo que asegura la terminación.
Si uno de los datos es cero, el resultado es el valor absoluto del otro.
El \textbf{algoritmo extendido} expresa ese divisor en función de los
datos iniciales, sustituyendo hacia atrás en las divisiones anteriores,
para obtener los coeficientes de Bézout.

\begin{ejemplo}
$84=2\cdot30+24$ y $30=24+6$, luego $\mcd(84,30)=6$.
Al sustituir hacia atrás,
$6=30-24=3\cdot30-84$.
\end{ejemplo}
